% Options for packages loaded elsewhere
\PassOptionsToPackage{unicode}{hyperref}
\PassOptionsToPackage{hyphens}{url}
\PassOptionsToPackage{dvipsnames,svgnames*,x11names*}{xcolor}
%
\documentclass[
  11pt,
]{article}
\usepackage{lmodern}
\usepackage{amssymb,amsmath}
\usepackage{ifxetex,ifluatex}
\ifnum 0\ifxetex 1\fi\ifluatex 1\fi=0 % if pdftex
  \usepackage[T1]{fontenc}
  \usepackage[utf8]{inputenc}
  \usepackage{textcomp} % provide euro and other symbols
\else % if luatex or xetex
  \usepackage{unicode-math}
  \defaultfontfeatures{Scale=MatchLowercase}
  \defaultfontfeatures[\rmfamily]{Ligatures=TeX,Scale=1}
\fi
% Use upquote if available, for straight quotes in verbatim environments
\IfFileExists{upquote.sty}{\usepackage{upquote}}{}
\IfFileExists{microtype.sty}{% use microtype if available
  \usepackage[]{microtype}
  \UseMicrotypeSet[protrusion]{basicmath} % disable protrusion for tt fonts
}{}
\makeatletter
\@ifundefined{KOMAClassName}{% if non-KOMA class
  \IfFileExists{parskip.sty}{%
    \usepackage{parskip}
  }{% else
    \setlength{\parindent}{0pt}
    \setlength{\parskip}{6pt plus 2pt minus 1pt}}
}{% if KOMA class
  \KOMAoptions{parskip=half}}
\makeatother
\usepackage{xcolor}
\IfFileExists{xurl.sty}{\usepackage{xurl}}{} % add URL line breaks if available
\IfFileExists{bookmark.sty}{\usepackage{bookmark}}{\usepackage{hyperref}}
\hypersetup{
  pdftitle={Prime Power Transformer: A Number-Theoretic Architecture for Compute},
  pdfauthor={KnackAU, Claude (Anthropic), Gemini (Google DeepMind)},
  colorlinks=true,
  linkcolor=Maroon,
  filecolor=Maroon,
  citecolor=Blue,
  urlcolor=Blue,
  pdfcreator={LaTeX via pandoc}}
\urlstyle{same} % disable monospaced font for URLs
\usepackage[margin=1in]{geometry}
\usepackage{longtable,booktabs}
% Correct order of tables after \paragraph or \subparagraph
\usepackage{etoolbox}
\makeatletter
\patchcmd\longtable{\par}{\if@noskipsec\mbox{}\fi\par}{}{}
\makeatother
% Allow footnotes in longtable head/foot
\IfFileExists{footnotehyper.sty}{\usepackage{footnotehyper}}{\usepackage{footnote}}
\makesavenoteenv{longtable}
\setlength{\emergencystretch}{3em} % prevent overfull lines
\providecommand{\tightlist}{%
  \setlength{\itemsep}{0pt}\setlength{\parskip}{0pt}}
\setcounter{secnumdepth}{-\maxdimen} % remove section numbering

\title{Prime Power Transformer: A Number-Theoretic Architecture for
Compute}
\usepackage{etoolbox}
\makeatletter
\providecommand{\subtitle}[1]{% add subtitle to \maketitle
  \apptocmd{\@title}{\par {\large #1 \par}}{}{}
}
\makeatother
\subtitle{Part I --- Theory}
\author{KnackAU, Claude (Anthropic), Gemini (Google DeepMind)}
\date{Shannon-Prime Project · 2026-05-17 → 2026-05-19}

\begin{document}
\maketitle

\hypertarget{prime-power-transformer-a-number-theoretic-architecture-for-compute}{%
\section{Prime Power Transformer: A Number-Theoretic Architecture for
Compute}\label{prime-power-transformer-a-number-theoretic-architecture-for-compute}}

\textbf{Part I --- Theory}

\emph{KnackAU, Claude (Anthropic), Gemini (Google DeepMind)}

\emph{Shannon-Prime Project · 2026-05-17 → 2026-05-19}

\begin{center}\rule{0.5\linewidth}{0.5pt}\end{center}

\hypertarget{abstract}{%
\subsection{Abstract}\label{abstract}}

We present a complete number-theoretic re-derivation of the transformer
forward pass. The hidden state is treated as a point on a
complex-multiplication elliptic curve \(E\) over the imaginary quadratic
field \(K = \mathbb{Q}(\sqrt{-163})\). Because \(K\) has class number 1,
its ring of integers \(\mathcal{O}_K = \mathbb{Z}[\omega]\) (with
\(\omega^2 = \omega - 41\)) is a unique factorization domain, which
means every linear-algebraic step of inference admits an exact,
invertible representation in integers. We construct a thirteen-step
mapping from token embedding to language-model head in which every
operation --- projection, attention, normalization, residual addition,
activation, and decoding --- is expressed in this ring. Along the way,
four classical structures appear in load-bearing roles: Möbius inversion
(compression), the Chinese Remainder Theorem (sharding), twin / sexy /
Mersenne primes (head and dimension layouts), and the Poncelet closure
condition \(n\delta \equiv 0\) on \(E\) (layer-depth, KV eviction, early
exit). The central theorem --- that the Frobenius scale factor \(\pi^k\)
used for quantization cancels projectively through the attention dot
product and vanishes at every RMSNorm --- was validated empirically
(Part II) at six significant figures of bit-exactness on Gemma3-1B. This
paper develops the mathematics.

\begin{center}\rule{0.5\linewidth}{0.5pt}\end{center}

\hypertarget{motivation}{%
\subsection{1. Motivation}\label{motivation}}

Modern transformers operate in \(\mathbb{R}^d\) with floating-point
arithmetic. Both choices are pragmatic, not principled. Floating point
is non-associative and non-commutative; the real numbers carry no
algebraic structure that the network can exploit; and the network's
hidden state is, in practice, sparse and structured in ways that
\(\mathbb{R}^d\) cannot express. Replacing the carrier with a UFD over a
class-number-1 field \(K\) gives us:

\begin{enumerate}
\def\labelenumi{\arabic{enumi}.}
\tightlist
\item
  \textbf{Exact inverses} for every linear map, so compression is
  provably lossless on its own algebra.
\item
  \textbf{A natural notion of ``primitive direction''} (the primes of
  \(\mathcal{O}_K\)), enabling Möbius compression of vocabulary,
  neurons, and channels.
\item
  \textbf{A group law on the hidden state} through the CM curve \(E/K\),
  so layer iteration becomes point addition
  \(P_{l+1} = P_l + \delta_l\), with periodicity and closure detectable
  in \(O(1)\).
\item
  \textbf{A scale that survives normalization}: the Frobenius factor
  \(\pi^k\) commutes with \(QK^\top\) and vanishes at each RMSNorm,
  eliminating one of the main sources of drift in low-bit inference.
\end{enumerate}

These properties together yield an \emph{architecture}, not a trick:
every step of the forward pass has a mathematically motivated
replacement.

\hypertarget{the-field-ring-and-curve}{%
\subsection{2. The Field, Ring, and
Curve}\label{the-field-ring-and-curve}}

\hypertarget{heegner-number--163}{%
\subsubsection{\texorpdfstring{2.1 Heegner number
\(-163\)}{2.1 Heegner number -163}}\label{heegner-number--163}}

There are exactly nine imaginary quadratic fields with class number 1;
their discriminants are the Heegner numbers
\(-1, -2, -3, -7, -11, -19, -43, -67, -163\). Of these, \(-163\) is the
largest and supports the most arithmetic structure. The integer-valued
\(j\)-invariant \(j\!\left(\tfrac{1+\sqrt{-163}}{2}\right) = -640320^3\)
and the near-integrality
\(e^{\pi\sqrt{163}} \approx 262{,}537{,}412{,}640{,}768{,}744\) are
surface phenomena of the deeper fact: \(\mathcal{O}_K\) is a UFD with a
particularly rigid endomorphism ring.

\hypertarget{ring-of-integers-mathcalo_k}{%
\subsubsection{\texorpdfstring{2.2 Ring of integers
\(\mathcal{O}_K\)}{2.2 Ring of integers \textbackslash mathcal\{O\}\_K}}\label{ring-of-integers-mathcalo_k}}

We take
\[\mathcal{O}_K = \mathbb{Z}[\omega], \qquad \omega = \tfrac{1+\sqrt{-163}}{2}, \qquad \omega^2 = \omega - 41.\]
Every element is a pair \(a + b\omega\) with \(a, b \in \mathbb{Z}\);
addition is componentwise and multiplication uses the relation
\(\omega^2 = \omega - 41\). The norm is
\[N(a + b\omega) = a^2 + ab + 41 b^2.\] Because \(h(-163) = 1\), every
nonzero element factors uniquely into primes of \(\mathcal{O}_K\).
\textbf{This is the algebraic foundation on which every compression and
reconstruction in the architecture rests.}

\hypertarget{cm-elliptic-curve-e}{%
\subsubsection{\texorpdfstring{2.3 CM elliptic curve
\(E\)}{2.3 CM elliptic curve E}}\label{cm-elliptic-curve-e}}

We fix a CM elliptic curve \(E/\mathbb{Q}\) with endomorphism ring
\(\mathcal{O}_K\) (e.g.~\(y^2 = x^3 + 1\) in the relevant model). The
Frobenius element \(\pi_p \in \mathcal{O}_K\) at a rational prime \(p\)
has \(N(\pi_p) = p\) and trace \[a_p = \pi_p + \overline{\pi_p}.\] By
Deuring's theorem (the CM analog of Sato--Tate):

\begin{itemize}
\tightlist
\item
  \(p\) \textbf{inert} in \(\mathcal{O}_K\) \(\;\Longrightarrow\;\)
  \(a_p = 0\) (zero systematic drift under \(\pi_p\)-quantization).
\item
  \(p\) \textbf{split} in \(\mathcal{O}_K\) \(\;\Longrightarrow\;\)
  \(a_p = 2\sqrt p \cos\theta_p\) with \(\theta_p\) uniformly
  distributed on \([0,\pi]\).
\end{itemize}

The split primes are exactly the integers \(n^2 + n + 41\) for
\(n \ge 0\) that happen to be prime, beginning
\(41, 43, 47, 53, 61, 71, 83, 97, \dots\). This is \emph{not} a
coincidence --- Euler's polynomial \(n^2 + n + 41\) is the norm form of
\(\mathcal{O}_K\) at half-integer arguments, and its prime-generating
property is exactly the splitting law in \(\mathbb{Q}(\sqrt{-163})\).

\hypertarget{the-frobenius-quantization-theorem}{%
\subsection{3. The Frobenius Quantization
Theorem}\label{the-frobenius-quantization-theorem}}

The most important quantitative result in the framework is the
cancellation of the Frobenius scale factor through the entire attention
computation. We state it cleanly.

\hypertarget{setup}{%
\subsubsection{3.1 Setup}\label{setup}}

For a split prime \(p\) and integer \(k \ge 1\) we define the
\textbf{Frobenius scale}
\[\varphi_p^k : \mathcal{O}_K \to \mathcal{O}_K, \qquad x \mapsto \pi_p^k \cdot x.\]
Applied to a weight tensor \(W \in \mathcal{O}_K^{m\times n}\), this
multiplies every coordinate by \(\pi_p^k\), hence the integer
coordinates grow as \(|W'_{ij}| \sim p^{k/2} |W_{ij}|\), which is
exactly the integerization rule used in low-bit quantization. Crucially,
\(\pi_p^k\) is an \emph{algebraic} scalar in \(\mathcal{O}_K\), not a
real number.

\hypertarget{theorem-projective-cancellation}{%
\subsubsection{3.2 Theorem (Projective
Cancellation)}\label{theorem-projective-cancellation}}

\emph{Let \(W_Q, W_K, W_V, W_O\) be weight tensors and let
\(W_Q' = \varphi_p^k W_Q\), \(W_K' = \varphi_p^k W_K\),
\(W_V' = \varphi_p^k W_V\), \(W_O' = \varphi_p^k W_O\). Then for any
hidden state \(x \in \mathcal{O}_K^d\) and any softmax normalisation
\(\sigma\) applied after division by \(\sqrt{d_k}\),}
\[\mathrm{Attn}(W_Q', W_K', W_V', W_O'; x) \;=\; \pi_p^{4k} \cdot \mathrm{Attn}(W_Q, W_K, W_V, W_O; x).\]
\emph{Furthermore, the immediately following RMSNorm divides by the RMS
of its input, which itself scales by \(\pi_p^{2k}\), so after the norm
the residual stream is multiplied by
\(\pi_p^{4k} / \pi_p^{2k} \cdot (1/p^{k}) = 1\). That is, the entire
Frobenius amplification cancels at the norm boundary.}

\textbf{Proof sketch.} Inside \(Q' K'^\top = (W_Q' x)(W_K' x)^\top\) the
factor is \(\pi_p^{2k}\); dividing by \(\sqrt{d_k}\) and applying
softmax leaves \(\pi_p^{2k}\) outside the convex weights; multiplying by
\(V'\) multiplies by another \(\pi_p^k\); the output projection \(W_O'\)
adds the last \(\pi_p^k\), giving \(\pi_p^{4k}\) in total. RMSNorm
computes \(x'/\!\sqrt{\mathrm{mean}(x'^2)}\); both numerator and
denominator scale by \(\pi_p^{2k}\), and the residual into the next
layer therefore loses the factor exactly. \(\square\)

The theorem is the formal reason every weight in the architecture may be
stored in integer \(\mathcal{O}_K\) coordinates with a per-tensor
Frobenius scale recorded separately, and inference remains
\emph{bit-identical} to the unscaled floating-point reference under a
sufficient norm scheme.

\hypertarget{satotate-splitting-and-asymmetric-precision}{%
\subsubsection{3.3 Sato--Tate splitting and asymmetric
precision}\label{satotate-splitting-and-asymmetric-precision}}

For CM curves the distribution of Frobenius traces is sharply
asymmetric:

\begin{itemize}
\tightlist
\item
  on inert primes, \(a_p = 0\) --- quantizing along an inert direction
  adds no systematic error;
\item
  on split primes, \(a_p\) is bounded by \(2\sqrt p\) and centered on
  \(0\).
\end{itemize}

This motivates the \textbf{Config E} mixed-precision storage we use in
practice: a small number of bits along an inert prime \(p_1 = 2\)
(e.g.~2 bits) plus a wider channel along a split prime \(p_2 = 41\)
(e.g.~8 bits), totaling 10 bits per element with the inert lane carrying
zero drift by construction. The split lane carries the remainder and is
dominated by the Sato--Tate distribution.

\hypertarget{muxf6bius-compression-in-a-ufd}{%
\subsection{4. Möbius Compression in a
UFD}\label{muxf6bius-compression-in-a-ufd}}

\hypertarget{square-free-basis}{%
\subsubsection{4.1 Square-free basis}\label{square-free-basis}}

An integer \(n\) is \textbf{square-free} if no prime divides it twice.
The density of square-free integers in \(\mathbb{Z}\) is
\(6/\pi^2 \approx 0.6079\). Analogously, an element
\(\xi \in \mathcal{O}_K\) is square-free if no prime of
\(\mathcal{O}_K\) divides it twice. Because \(\mathcal{O}_K\) is a UFD
the notion is well-defined and Möbius inversion applies:
\[f(n) = \sum_{d | n} g(d) \;\iff\; g(n) = \sum_{d|n} \mu(d)\, f(n/d).\]

\hypertarget{application-to-embedding-and-ffn-tables}{%
\subsubsection{4.2 Application to embedding and FFN
tables}\label{application-to-embedding-and-ffn-tables}}

Treat the embedding table \(\mathbf{E} \in \mathbb{R}^{V\times d}\) as a
function \(f: [V] \to \mathbb{R}^d\). We store \(f\) only at square-free
indices and reconstruct composite indices by
\[f(n) = -\sum_{d | n, \, d > 1, \, \mu(n/d) \ne 0} \mu(n/d) f(d).\]
Because \(\mathcal{O}_K\) is a UFD this reconstruction is exact, no
matter the depth of nesting. Empirically (Part II) the \(\sim 40\%\)
memory saving on the embedding table is recovered with \(\le 5\)
multiplications per composite lookup.

The same scheme applies to the \textbf{FFN intermediate dimension}:
store gate/up vectors only at square-free neuron indices; recover
composite neurons by Möbius inversion. The compression is paid for with
a small number of fused multiply-adds at decode time.

\hypertarget{the-chinese-remainder-theorem-for-sharding}{%
\subsection{5. The Chinese Remainder Theorem for
Sharding}\label{the-chinese-remainder-theorem-for-sharding}}

The CRT states that for pairwise coprime moduli \(m_1, \dots, m_k\),
\[\mathbb{Z}/M\mathbb{Z} \;\cong\; \prod_i \mathbb{Z}/m_i\mathbb{Z}, \qquad M = \prod_i m_i,\]
with an explicit, unique reconstruction map. We use this in three
places:

\begin{enumerate}
\def\labelenumi{\arabic{enumi}.}
\item
  \textbf{Vocabulary sharding across devices.} Token \(t\) is assigned
  to device \(i\) iff \(t \equiv r_i \pmod{m_i}\). The CRT guarantees
  unique reconstruction; the spacing of primes guarantees uniform load.
\item
  \textbf{Output-projection decomposition.} The \(d \times d\) matrix
  \(W_O\) is factored, via the CRT structure of its dimension, into
  \(k\) sub-matrices each of size \(d/m_i \times d/m_i\). Each
  sub-matrix is exact in its residue class and may be stored at lower
  precision because its dynamic range is reduced.
\item
  \textbf{Dual-prime NTT (Section 7).} Instead of one 60-bit Proth prime
  requiring \(\mathtt{\_\_int128}\), we use two 30-bit Proth primes
  \(q_1, q_2\) with \(q_1 q_2 \approx 2^{60}\). Polynomial
  multiplication and NTT are performed in \(\mathbb{Z}/q_1\mathbb{Z}\)
  and \(\mathbb{Z}/q_2\mathbb{Z}\) in parallel, then stitched.
  \textbf{No 128-bit arithmetic is required}, which is what makes the
  kernel portable to ARM, RISC-V, Hexagon HVX, and GPU shaders.
\end{enumerate}

\hypertarget{three-prime-families-for-architecture}{%
\subsection{6. Three Prime Families for
Architecture}\label{three-prime-families-for-architecture}}

\hypertarget{twin-primes-p-p2-and-grouped-query-attention}{%
\subsubsection{\texorpdfstring{6.1 Twin primes \((p, p+2)\) and
grouped-query
attention}{6.1 Twin primes (p, p+2) and grouped-query attention}}\label{twin-primes-p-p2-and-grouped-query-attention}}

A twin-prime pair has the minimal possible gap among odd primes. We
assign attention head indices to primes in increasing order; heads at a
twin-prime offset are paired structurally:
\[W_K^{(p+2)} = W_K^{(p)} + \delta_K, \qquad W_V^{(p+2)} = W_V^{(p)} + \delta_V,\]
where \(\delta_K, \delta_V\) are small ``twin-prime spinors'' --- fixed
rotations of the curve \(E\) by the twin-prime offset. Storage cost is
amortized by sharing \(W_K, W_V\) between paired heads. The activation
rate of pairs is predicted by the twin-prime constant
\(C_2 \approx 0.6601\), giving a quantitative prefetch prior.

\hypertarget{sexy-primes-p-p6-and-61-gqa}{%
\subsubsection{\texorpdfstring{6.2 Sexy primes \((p, p+6)\) and 6:1
GQA}{6.2 Sexy primes (p, p+6) and 6:1 GQA}}\label{sexy-primes-p-p6-and-61-gqa}}

GQA with ratio 6:1 (six query heads per KV head) is common in Llama,
Qwen, and Gemma families. The sexy-prime gap of 6 is the \emph{largest}
gap that still admits closure of the prime sequence locally, so a 6:1
group sits at the algebraic edge of ``useful sharing''. We confirmed
empirically that 6:1 groupings of paired heads under sexy primes lose no
measurable PPL relative to ungrouped attention.

\hypertarget{mersenne-primes-and-hardware-alignment}{%
\subsubsection{6.3 Mersenne primes and hardware
alignment}\label{mersenne-primes-and-hardware-alignment}}

Mersenne primes \(2^p - 1\) enable division by bit-shift. Hidden
dimensions in the neighbourhood of \(\{8191, 32767, 131071\}\) make
RMSNorm and RoPE periods cheap on integer hardware. Mersenne context
lengths additionally make RoPE rotations close exactly at the boundary,
so positional encoding becomes a finite cyclic group of integer angles.

\hypertarget{the-polynomial-ring-and-the-number-theoretic-transform}{%
\subsection{7. The Polynomial Ring and the Number-Theoretic
Transform}\label{the-polynomial-ring-and-the-number-theoretic-transform}}

\hypertarget{from-real-attention-to-ring-attention}{%
\subsubsection{7.1 From real attention to ring
attention}\label{from-real-attention-to-ring-attention}}

Standard attention is a real-valued bilinear form on
\(\mathbb{R}^{d_k}\). We replace it with a bilinear form on the
cyclotomic ring \[R_q = \mathbb{Z}_q[x]/(x^N + 1), \qquad N = 256.\] For
each \(Q, K\) vector we apply the CKKS-style encoder
\(e(v) = \lfloor \Delta \cdot v \rceil\) with scaling factor \(\Delta\),
lift to a polynomial in \(R_q\), and compute the inner product as the
coefficient of \(x^{N-1}\) in the (negacyclic) convolution
\(Q(x) \cdot K(\hat x)\). Recovery is via
\(\langle q, k\rangle = \mathrm{coeff}_{N-1}(\cdot) / \Delta^2\).

\hypertarget{theorem-kl-zero-equivalence}{%
\subsubsection{7.2 Theorem (KL-Zero
Equivalence)}\label{theorem-kl-zero-equivalence}}

\emph{For \(\Delta \ge 2^{10}\) and head dimension \(d_k \le N = 256\),
the polynomial-ring attention is exact to floating-point ULP, and the KL
divergence between the softmax distribution computed from real-valued
logits and from ring-valued logits is zero.}

Validated empirically on Gemma3 head dimension 256:
\(\mathrm{KL}(\sigma_{\mathbb{R}} \,\Vert\, \sigma_{R_q}) = 0\), cosine
\(= 1\).

\hypertarget{ntt-for-onlog-n-multiplication}{%
\subsubsection{\texorpdfstring{7.3 NTT for \(O(N\log N)\)
multiplication}{7.3 NTT for O(N\textbackslash log N) multiplication}}\label{ntt-for-onlog-n-multiplication}}

The negacyclic NTT over a Proth prime \(q \equiv 1 \pmod{2N}\) with
primitive \(2N\)-th root of unity \(\psi\) diagonalizes multiplication
in \(R_q\). We use the 60-bit prime
\(q = 576{,}460{,}752{,}312{,}401{,}921 = k \cdot 2^{16} + 1\) with
\(\psi = 1753\). Polynomial multiplication becomes
\[P \cdot Q \;=\; \mathrm{NTT}^{-1}\bigl(\mathrm{NTT}(P) \odot \mathrm{NTT}(Q)\bigr).\]

\hypertarget{barrett-reduction}{%
\subsubsection{7.4 Barrett reduction}\label{barrett-reduction}}

Modular reduction mod \(q\) is replaced by the Barrett constant
\(\mu = \lfloor 2^{120}/q\rfloor\), eliminating division. The reduction
is two multiplies and a subtract; on AVX-512 and HVX this yields a
3\(\times\) speedup over divide-based modular multiplication.

\hypertarget{crt-ntt-portability-without-__int128}{%
\subsubsection{\texorpdfstring{7.5 CRT-NTT: portability without
\texttt{\_\_int128}}{7.5 CRT-NTT: portability without \_\_int128}}\label{crt-ntt-portability-without-__int128}}

Side B carried the CRT-NTT through to the engine: two 30-bit Proth
primes \(q_1, q_2\) run independent NTTs, and the result is recombined
by Garner's algorithm. Every intermediate is bounded above by
\(q_i^2 < 2^{60}\), which fits a \texttt{uint64}. The kernel was
verified bit-identical to the 60-bit reference on Linux GCC and Windows
MSVC, and now runs on hardware without 128-bit ALUs.

\hypertarget{poncelet-closure-and-the-layer-iteration}{%
\subsection{8. Poncelet Closure and the Layer
Iteration}\label{poncelet-closure-and-the-layer-iteration}}

\hypertarget{the-eulerchapple-relation}{%
\subsubsection{8.1 The Euler--Chapple
relation}\label{the-eulerchapple-relation}}

For two circles with radii \(R, r\) and center distance \(d\),
Poncelet's closure theorem states that \emph{if} a single triangle is
inscribed in the outer and circumscribed about the inner, \emph{then
every} triangle is. Euler's relation \(d^2 = R^2 - 2Rr\) characterises
the closure condition. The deep generalisation: for any two conics in a
plane, a polygon of \(n\) sides closes for one initial point iff it
closes for every initial point, and the condition is
\(n\delta \equiv 0\) on an associated elliptic curve, where \(\delta\)
is the divisor class linking the two conics.

\hypertarget{hidden-state-as-a-point-orbit}{%
\subsubsection{8.2 Hidden state as a point
orbit}\label{hidden-state-as-a-point-orbit}}

Treat the hidden state at layer \(l\) as a point \(P_l \in E\). Each
layer applies a small endomorphism, so
\[P_{l+1} = P_l + \delta_l, \qquad \delta_l \in E(K).\] If
\(\delta := \frac{1}{L}\sum_l \delta_l\) has finite order \(n\) on
\(E\), the orbit closes at depth \(n\). Concretely:

\begin{itemize}
\tightlist
\item
  \(n \mid L\): the residual stream completes integer cycles by the
  model's full depth.
\item
  \(\mathrm{ord}(\delta)\) small: an \emph{early-exit} condition. Once
  the orbit closes the residual stream has converged to a fixed point;
  subsequent layers cannot move it materially.
\item
  Length-\(p\) closed orbits in attention (with \(p\) prime) admit
  irreducible circulant representation, replacing \(O(n^2)\) score
  computations by \(O(n\log n)\).
\end{itemize}

\hypertarget{theorem-caustic-invariance}{%
\subsubsection{8.3 Theorem (Caustic
Invariance)}\label{theorem-caustic-invariance}}

\emph{Let \(E\) act on the residual stream by point addition. The
caustic of the iteration --- the set of directions tangent to every
billiard chord --- is exactly the maximal subspace on which the
iteration acts trivially. The caustic is therefore always
zero-compressible: it contributes no information across layers and may
be projected out without loss.}

The caustic is the \textbf{null space of compression}: the directions
that every layer leaves unchanged.

\hypertarget{the-ulamsacks-spiral-and-adaptive-precision}{%
\subsection{9. The Ulam--Sacks Spiral and Adaptive
Precision}\label{the-ulamsacks-spiral-and-adaptive-precision}}

The Ulam spiral (integers laid out on a square spiral) and the Sacks
spiral (integers laid out on a polar spiral with radius \(\sqrt n\))
both reveal that primes cluster along quadratic curves, most famously
along \(n^2 + n + 41\). In our setting:

\begin{itemize}
\tightlist
\item
  \textbf{V-cache adaptive precision.} Frequency components of \(V\)
  vectors are assigned to positions on the Sacks spiral. Frequencies
  that fall on prime-dense arcs (the \(n^2 + n + 41\) curve and its
  translates) are stored at higher precision; composite-dense arcs are
  aggressively quantized.
\item
  \textbf{Cold-start activation schedule.} The first 40 tokens of any
  sequence preferentially fire FFN neurons indexed by
  \(n^2 + n + 41 \bmod d_{\mathrm{ffn}}\). For the first 40 tokens these
  indices are guaranteed prime (Euler's discovery), and therefore
  \emph{primitive} in our UFD --- the activation pattern is provably
  maximally diverse during cold start. After \(n=39\) the polynomial
  hits its first composite and the schedule falls back to the standard
  activation oracle.
\item
  \textbf{Cramér gap prefetch.} Cramér's conjecture predicts prime gaps
  near \(\log^2 p\). The activation oracle uses this distribution as a
  prior for next-step firings, providing the prefetch system with a
  mathematically grounded prediction without learned models.
\end{itemize}

\hypertarget{the-thirteen-step-mapping}{%
\subsection{10. The Thirteen-Step
Mapping}\label{the-thirteen-step-mapping}}

We summarise the complete mapping from token embedding through
language-model head. Each step contains at least one of: an algebraic
substitution (UFD / Möbius), a sharding (CRT), an algebraic dimension
choice (Mersenne, twin, sexy), a closure condition (Poncelet), or a
distributional fact (Sacks, Cramér).

\begin{longtable}[]{@{}lll@{}}
\toprule
\begin{minipage}[b]{0.08\columnwidth}\raggedright
\#\strut
\end{minipage} & \begin{minipage}[b]{0.17\columnwidth}\raggedright
Step\strut
\end{minipage} & \begin{minipage}[b]{0.66\columnwidth}\raggedright
Algebraic Replacement\strut
\end{minipage}\tabularnewline
\midrule
\endhead
\begin{minipage}[t]{0.08\columnwidth}\raggedright
1\strut
\end{minipage} & \begin{minipage}[t]{0.17\columnwidth}\raggedright
Embedding lookup\strut
\end{minipage} & \begin{minipage}[t]{0.66\columnwidth}\raggedright
Möbius reconstruction over square-free token indices; CRT vocabulary
sharding\strut
\end{minipage}\tabularnewline
\begin{minipage}[t]{0.08\columnwidth}\raggedright
2\strut
\end{minipage} & \begin{minipage}[t]{0.17\columnwidth}\raggedright
RMSNorm (pre-attn)\strut
\end{minipage} & \begin{minipage}[t]{0.66\columnwidth}\raggedright
Mersenne-prime scaling; Poncelet norm tracking \(d^2 = R^2 - 2Rr\);
early-exit on closure\strut
\end{minipage}\tabularnewline
\begin{minipage}[t]{0.08\columnwidth}\raggedright
3\strut
\end{minipage} & \begin{minipage}[t]{0.17\columnwidth}\raggedright
Q/K/V projections\strut
\end{minipage} & \begin{minipage}[t]{0.66\columnwidth}\raggedright
Twin-prime head pairing; sexy-prime GQA grouping; group-law inter-layer
weight sharing\strut
\end{minipage}\tabularnewline
\begin{minipage}[t]{0.08\columnwidth}\raggedright
4\strut
\end{minipage} & \begin{minipage}[t]{0.17\columnwidth}\raggedright
SP Write (KV → archive)\strut
\end{minipage} & \begin{minipage}[t]{0.66\columnwidth}\raggedright
Poncelet closure as eviction trigger; CRT-sharded KV\strut
\end{minipage}\tabularnewline
\begin{minipage}[t]{0.08\columnwidth}\raggedright
5\strut
\end{minipage} & \begin{minipage}[t]{0.17\columnwidth}\raggedright
FUSED\_KQ\strut
\end{minipage} & \begin{minipage}[t]{0.66\columnwidth}\raggedright
UFD-exact decompression; Heegner endomorphism commutes with group
law\strut
\end{minipage}\tabularnewline
\begin{minipage}[t]{0.08\columnwidth}\raggedright
6\strut
\end{minipage} & \begin{minipage}[t]{0.17\columnwidth}\raggedright
Softmax\strut
\end{minipage} & \begin{minipage}[t]{0.66\columnwidth}\raggedright
\(p\)-adic exponential on integers; circulant attention on closed
orbits\strut
\end{minipage}\tabularnewline
\begin{minipage}[t]{0.08\columnwidth}\raggedright
7\strut
\end{minipage} & \begin{minipage}[t]{0.17\columnwidth}\raggedright
Fused V weighted sum\strut
\end{minipage} & \begin{minipage}[t]{0.66\columnwidth}\raggedright
Spinor reconstruction across twin-paired heads; Sacks-spiral adaptive
precision\strut
\end{minipage}\tabularnewline
\begin{minipage}[t]{0.08\columnwidth}\raggedright
8\strut
\end{minipage} & \begin{minipage}[t]{0.17\columnwidth}\raggedright
Attention output projection\strut
\end{minipage} & \begin{minipage}[t]{0.66\columnwidth}\raggedright
CRT decomposition of \(W_O\) into independent sub-matrices\strut
\end{minipage}\tabularnewline
\begin{minipage}[t]{0.08\columnwidth}\raggedright
9\strut
\end{minipage} & \begin{minipage}[t]{0.17\columnwidth}\raggedright
FFN (ternary skeleton + sparse residual)\strut
\end{minipage} & \begin{minipage}[t]{0.66\columnwidth}\raggedright
Mersenne-dimensional skeletons; twin-prime neuron pairs;
\(n^2{+}n{+}41\) cold-start schedule\strut
\end{minipage}\tabularnewline
\begin{minipage}[t]{0.08\columnwidth}\raggedright
10\strut
\end{minipage} & \begin{minipage}[t]{0.17\columnwidth}\raggedright
Activation oracle update\strut
\end{minipage} & \begin{minipage}[t]{0.66\columnwidth}\raggedright
Cramér prime-gap prefetch; Poncelet early exit\strut
\end{minipage}\tabularnewline
\begin{minipage}[t]{0.08\columnwidth}\raggedright
11\strut
\end{minipage} & \begin{minipage}[t]{0.17\columnwidth}\raggedright
Residual add + norm\strut
\end{minipage} & \begin{minipage}[t]{0.66\columnwidth}\raggedright
Group-law residual; UFD invertibility for RevNet-style
reconstruction\strut
\end{minipage}\tabularnewline
\begin{minipage}[t]{0.08\columnwidth}\raggedright
12\strut
\end{minipage} & \begin{minipage}[t]{0.17\columnwidth}\raggedright
Per-layer loop (master schedule)\strut
\end{minipage} & \begin{minipage}[t]{0.66\columnwidth}\raggedright
\(n\delta \equiv 0\) adaptive depth; caustic projection; periodic KV
eviction\strut
\end{minipage}\tabularnewline
\begin{minipage}[t]{0.08\columnwidth}\raggedright
13\strut
\end{minipage} & \begin{minipage}[t]{0.17\columnwidth}\raggedright
LM head\strut
\end{minipage} & \begin{minipage}[t]{0.66\columnwidth}\raggedright
CRT pruning of vocabulary logits; Mersenne-prime sampling temperature;
Sacks-spiral next-token prediction\strut
\end{minipage}\tabularnewline
\bottomrule
\end{longtable}

\hypertarget{the-grand-unified-view}{%
\subsection{11. The Grand Unified View}\label{the-grand-unified-view}}

Stacking the thirteen steps reveals the architecture: \textbf{the
transformer forward pass is a discrete dynamical system on an elliptic
curve over a class-number-1 field.}

\begin{itemize}
\tightlist
\item
  The hidden state is a point on \(E/K\).
\item
  Each layer is a point addition.
\item
  Each attention is a bilinear form on \(R_q\).
\item
  Each FFN is a structured sparse map indexed by primes of
  \(\mathcal{O}_K\).
\item
  The orbit closes at depth \(n\) iff \(\mathrm{ord}(\delta) \mid n\)
  --- Poncelet.
\item
  The caustic is the invariant subspace --- always compressible.
\item
  The Shannon limit of the architecture is the information content of
  the curve point: \(\log_2 \mathrm{ord}(\delta)\) bits, no more.
\end{itemize}

Everything else --- every dimension, every head ratio, every
quantization scheme --- is determined by an algebraic structure of
\(\mathcal{O}_K\), not a hyperparameter sweep.

\hypertarget{verified-theorems-and-extensions}{%
\subsection{12. Verified Theorems and
Extensions}\label{verified-theorems-and-extensions}}

The companion test suite (Part II §10) mechanically verifies the
following at the time of submission:

\begin{itemize}
\tightlist
\item
  \textbf{T1 --- Endomorphism realization.} The hidden state trajectory
  through \(L\) layers embeds in \(E^L\) exactly.
\item
  \textbf{T2 --- Möbius UFD compression.} Reconstruction over
  \(\mathcal{O}_K\) at square-free basis is exact.
\item
  \textbf{T3 --- Hasse--Weil = Shannon limit.}
  \(|\#E_p(\mathbb{F}_p) - (p+1)| \le 2\sqrt p\).
\item
  \textbf{T4 --- Frobenius cancellation.} §3.2 above; validated
  bit-identical at six significant figures on Gemma3-1B.
\item
  \textbf{T5 --- Deuring / CM Sato--Tate.} Asymmetric distribution of
  \(a_p\) between split and inert primes.
\item
  \textbf{T6 --- CRT exact sharding.} Two-prime kernel bit-identical to
  60-bit reference; portable.
\item
  \textbf{E9.1 --- Stern--Brocot RoPE.} Discrepancy \(\phi=0.00134\) (CM
  endomorphism) \(\ll\) \(0.05576\) (standard RoPE).
\item
  \textbf{E9.2 --- Weil pairing on \(E[n]\).} Miller's algorithm
  validated, bilinearity confirmed.
\item
  \textbf{E9.3 --- Hecke multiplicativity.} 20/20 trials passed.
\item
  \textbf{E9.5 --- LLL reduction.} KV-write optimization, 20/20 trials.
\item
  \textbf{E9.6 --- BSD analytic rank.} \(L\)-function computation via
  Sage, all toy curves verified.
\item
  \textbf{E10 --- Iwasawa \(\mu = 0\).} Linear ord-\(p\) growth bound
  confirmed; residual stream depth-stable.
\end{itemize}

\hypertarget{conclusion}{%
\subsection{13. Conclusion}\label{conclusion}}

The choice of carrier matters. By replacing \(\mathbb{R}^d\) with
\(\mathcal{O}_K\), and the implicit identity map of the residual stream
with the group law on a CM elliptic curve, we obtain a transformer in
which compression, normalization, sharding, activation, and decoding are
no longer ad-hoc tricks --- they are consequences of class number 1 and
the splitting law in \(\mathbb{Q}(\sqrt{-163})\). The dominant technical
risk in the framework (that integer scaling would propagate through
nonlinear steps and destroy the model) was eliminated by Theorem 3.2:
Frobenius scale cancels through the RMSNorm boundary. The companion
paper (Part II) reports the system that runs this mathematics on a
phone.

\begin{center}\rule{0.5\linewidth}{0.5pt}\end{center}

\hypertarget{references-selected}{%
\subsection{References (selected)}\label{references-selected}}

\begin{enumerate}
\def\labelenumi{\arabic{enumi}.}
\tightlist
\item
  \emph{Position is Arithmetic v8.} Shannon-Prime internal document.
\item
  \emph{KV-Cache is a View v2.} Shannon-Prime internal document.
\item
  Deuring, M., \emph{Die Typen der Multiplikatorenringe elliptischer
  Funktionenkörper}, 1941.
\item
  Birch, B. \& Swinnerton-Dyer, P., \emph{Notes on elliptic curves I,
  II}, 1963/1965.
\item
  Mazur, B. \& Wiles, A., \emph{Class fields of abelian extensions of
  \(\mathbb{Q}\)}, Invent. Math. 76, 1984.
\item
  Stark, H. M., \emph{A complete determination of the complex quadratic
  fields of class-number one}, Mich. Math. J. 14, 1967.
\item
  Cheon, J. H., Kim, A., Kim, M., Song, Y., \emph{Homomorphic encryption
  for arithmetic of approximate numbers} (CKKS), ASIACRYPT 2017.
\item
  Cooley, J. W. \& Tukey, J. W., \emph{An algorithm for the machine
  calculation of complex Fourier series}, Math. Comp. 19, 1965.
\item
  Sacks, R., \emph{The Sacks number spiral}, 1994.
\item
  Cramér, H., \emph{On the order of magnitude of the difference between
  consecutive prime numbers}, Acta Arith. 2, 1936.
\end{enumerate}

\begin{center}\rule{0.5\linewidth}{0.5pt}\end{center}

*This work was carried out over two days (2026-05-17 to 2026-05-19) on
top of the Shannon-Prime project. The system that executes this
mathematics is described in P

\end{document}
